\documentclass[../../main.tex]{subfiles}
\begin{document}

{\bf[解]}

\begin{figure}[htb]
  \centering
  \begin{tikzpicture}[scale=3.8, >=stealth, every node/.style={font=\small}]

    % --- 1. 定数・パラメータ定義 ---
    % 視認性の良い t の値 (0 < t < 1 の範囲で t = 0.55 を採用)
    \def\tVal{0.55}
    \pgfmathsetmacro{\xQ}{\tVal}
    \pgfmathsetmacro{\yQ}{\tVal*\tVal}                       % t^2 = 0.3025
    \pgfmathsetmacro{\sVal}{\tVal - \tVal*\tVal}             % s = t - t^2 = 0.2475
    \pgfmathsetmacro{\xP}{\xQ - 2.0*\tVal*\sVal}             % X = t - 2t(t-t^2) ≈ 0.2778
    \pgfmathsetmacro{\yP}{\yQ + \sVal}                       % Y = t^2 + (t-t^2) = t = 0.5500

    \coordinate (O) at (0, 0);
    \coordinate (Q) at (\xQ, \yQ);
    \coordinate (P) at (\xP, \yP);

    % --- 2. 座標軸 ---
    \draw[->, thick] (-0.3, 0) -- (1.3, 0) node[right] {$x$};
    \draw[->, thick] (0, -0.2) -- (0, 1.35) node[above] {$y$};
    \node[below left=2pt] at (O) {$O$};

    % --- 4. 放物線 y = x^2 ---
    \draw[very thick, black, domain=-0.25:1.12, samples=80]
    plot (\x, {\x*\x});
    \node[black, font=\footnotesize, above right=6pt] at (0.95, 0.95*0.95) {$y = x^2$};

    % --- 5. 接線 (傾き 2t) ---
    % 接線: y - yQ = 2t(x - xQ)
    \draw[thick, gray!70, domain=0.2:0.9]
    plot (\x, {\yQ + 2.0*\tVal*(\x - \xQ)});
    % \node[gray!70!black, font=\scriptsize, below right] at (0.85, {\yQ + 2.0*\tVal*(0.85 - \xQ)}) {接線};

    % --- 6. 法線および法線ベクトル \vec{l} ---
    % 法線全体 (直線)
    \draw[thick, blue!80!black]
    ($ (Q)!-0.35!(P) $) -- ($ (Q)!1.45!(P) $) node[above left, font=\footnotesize] {法線};

    % 線分 PQ (太線で強調)
    \draw[line width=1.8pt, orange!90!black] (Q) -- (P);

    % 点 Q を始点とする法線ベクトル \vec{l}
    \coordinate (VecL) at ($ (Q) + 0.35*(-2.0*\tVal, 1.0) $);
    \draw[->, very thick, red!85!black] (Q) -- (VecL)
    node[below left=1pt, font=\scriptsize] {$\vec{l} =
      \begin{pmatrix} -2t \\ 1
    \end{pmatrix}$};

    % 直角マーク (接線 ⊥ 法線)
    % 接線方向単位ベクトル (1, 2t)/sqrt(1+4t^2), 法線方向 (-2t, 1)/sqrt(1+4t^2)
    \pgfmathsetmacro{\lenNorm}{sqrt(1.0 + 4.0*\tVal*\tVal)}
    \pgfmathsetmacro{\dxT}{0.06 / \lenNorm}
    \pgfmathsetmacro{\dyT}{0.06 * 2.0*\tVal / \lenNorm}
    \pgfmathsetmacro{\dxN}{-0.06 * 2.0*\tVal / \lenNorm}
    \pgfmathsetmacro{\dyN}{0.06 / \lenNorm}
    \draw[thin, gray!70]
    ($ (Q) + (\dxT, \dyT) $) --
    ($ (Q) + (\dxT + \dxN, \dyT + \dyN) $) --
    ($ (Q) + (\dxN, \dyN) $);

    % --- 7. 座標補助線 ---
    \draw[dashed, gray!80!black] (\xQ, 0) -- (Q) -- (0, \yQ);
    \fill (\xQ, 0) circle (0.4pt) node[below=2pt, font=\scriptsize] {$t$};
    \fill (0, \yQ) circle (0.4pt) node[left=2pt, font=\scriptsize] {$t^2$};

    % --- 8. 主要点のプロットとラベル ---
    \fill[black] (Q) circle (0.7pt) node[below right=1pt] {$Q(t, t^2)$};
    \fill[orange!90!black] (P) circle (0.8pt) node[above right=1pt, orange!90!black] {$P(X, Y)$};

    % 長さ |PQ| の寸法表示
    \draw[<->, orange!90!black, thin]
    ($ (Q) + (0.05, 0.03) $) -- ($ (P) + (0.05, 0.03) $)
    node[midway, right=1pt, font=\tiny] {$|\overrightarrow{PQ}| = (t-t^2)\sqrt{1+4t^2}$};

    % 条件 y >= x^2 の領域注記
    \node[gray!60!black, font=\scriptsize] at (0.12, 0.95) {$y \ge x^2$};

  \end{tikzpicture}
  \caption{放物線 $y = x^2$ 上の点 $Q(t, t^2)$ における法線と動点 $P$}
  \label{fig:parabola_normal_vector_P}
\end{figure}

$(x^2)'=2x$だから，$Q$での法線方向のベクトル$\vec{l}$は，
\begin{align}
  \vec{l}=
  \begin{pmatrix} -2t \\ 1
  \end{pmatrix} \label{eq:1}
\end{align}
と書ける．故に法線上の点$P(X,Y)$は
\begin{align}
  \begin{pmatrix} X \\ Y
  \end{pmatrix}=
  \begin{pmatrix} t \\ t^2
  \end{pmatrix}+s
  \begin{pmatrix} -2t \\ 1
  \end{pmatrix} \label{eq:2}
\end{align}
と書ける．
この時
\begin{align}
  |\overrightarrow{PQ}| = |s|\sqrt{4t^2+1} \label{eq:3}
\end{align}
である．従って(ハ)の条件から，
\begin{align}
  &\left|s
  \begin{pmatrix} -2t \\ 1
  \end{pmatrix}\right|=(t-t^2)\sqrt{1+4t^2} \\
  \Longleftrightarrow &s=\pm(t-t^2) \label{eq:4}
\end{align}
である．ただし条件(ロ)から$s>0$だから絶対値を外した．

\cref{eq:4}を\cref{eq:2}に代入して，$P$は
\begin{align}
  \begin{pmatrix} X \\ Y
  \end{pmatrix}
  =
  \begin{pmatrix} t \\ t^2
  \end{pmatrix}+(t-t^2)
  \begin{pmatrix} -2t \\ 1
  \end{pmatrix}
\end{align}
で表されるパラメータ曲線のうち$y\ge x^2$の部分を動く．
$t$を消去して
\begin{align}
  X=2Y^3-2Y^2+Y \quad (0\le Y\le 1)
\end{align}
である．グラフの概形は下図のようになる．

\begin{figure}[htb]
  \centering
  \begin{tikzpicture}[scale=3.5, >=stealth, every node/.style={font=\small}]

    % --- 1. 求める面積 S の斜線ハッチング領域 (0 <= y <= 1) ---
    % x の範囲: 左側が x = f(y) = 2y^3 - 2y^2 + y, 右側が x = sqrt(y)
    \fill[pattern=north east lines, pattern color=blue!60, opacity=0.7]
    plot[domain=0:1, samples=60, variable=\y] ({2.0*\y*\y*\y - 2.0*\y*\y + \y}, {\y})
    -- plot[domain=1:0, samples=60, variable=\y] ({sqrt(\y)}, {\y})
    -- cycle;

    % --- 2. 座標軸 ---
    \draw[->, thick] (-0.2, 0) -- (1.35, 0) node[right] {$x$};
    \draw[->, thick] (0, -0.2) -- (0, 1.45) node[above] {$y$};
    \node[below left=2pt] at (0, 0) {$O$};

    % --- 3. 曲線 y = x^2 (x = \sqrt{y}) ---
    \draw[very thick, black, domain=-0.2:1.15, samples=80]
    plot (\x, {\x*\x});
    \node[black, font=\footnotesize, above right] at (0.8, 1.18) {$y = x^2$};

    % --- 4. 曲線 x = f(y) = 2y^3 - 2y^2 + y ---
    \draw[very thick, blue!85!black, domain=0:1.08, samples=80, variable=\y]
    plot ({2.0*\y*\y*\y - 2.0*\y*\y + \y}, {\y});

    % --- 5. 点 (1, 1) と座標補助線 ---
    \draw[dashed, gray!70] (1, 0) -- (1, 1) -- (0, 1);
    \fill (1, 1) circle (0.6pt);
    \node[below=2pt, font=\scriptsize] at (1, 0) {$1$};
    \node[left=2pt, font=\scriptsize] at (0, 1) {$1$};

    % --- 6. ラベル・引き出し線 ---
    % 曲線 f(y) への引き出し線（手書き図の配置に準拠）
    \draw[<-, thin, blue!85!black] (0.28, 0.42) to[bend right=20] (0.42, 0.15)
    node[right, font=\footnotesize, blue!90!black] {$x = f(y)$};

  \end{tikzpicture}
  \caption{放物線 $y = x^2$ と曲線 $x = f(y)$，および囲まれた領域の面積 $S$（斜線部）}
  \label{fig:parabola_normal_locus_area}
\end{figure}

求める面積$S$は上図の斜線部で，
\begin{align}
  S
  &= \int_{0}^{1} \left(\sqrt{y}-f(y)\right) dy \\
  &= \left[\dfrac{2}{3}y^{3/2}-\dfrac{1}{2}y^4+\dfrac{2}{3}y^3-\dfrac{1}{2}y^2\right]_{0}^{1} \\
  &= \dfrac{1}{3}
\end{align}
である．$\cdots$(答)

\newpage
\end{document}
